\documentclass{report}
\usepackage{graphicx} % Required for inserting images
\usepackage{amsmath}
\usepackage{amsthm}
\usepackage{amssymb}
\usepackage[a4paper, left=2cm, right=2cm, top=2cm, bottom=2cm]{geometry}
\usepackage[table]{xcolor} % Necessario per colorare righe/colonne
\usepackage{pifont}
\definecolor{SALZO}{rgb}{0.9,0.9,0.9}
\definecolor{lightblue}{rgb}{0.78, 0.94, 1.00}
\setcounter{secnumdepth}{1}
\usepackage{tikz} % Per disegnare il diagramma
\usepackage{hyperref}
\usepackage{bm}
\usepackage{cancel}
\usepackage[italian]{babel}
\usepackage{authblk}
\usepackage{pgfplots} % Pacchetto per i grafici
\pgfplotsset{compat=1.18} % Compatibilità con la versione
\usepackage{subcaption}
\usepackage{footnote}
\usepackage{listings}
\usepackage{tcolorbox}
\tcbuselibrary{listings}
\newtheoremstyle{nostyle}
  {0pt} % Spazio sopra
  {0pt} % Spazio sotto
  {\normalfont} % Corpo normale (non corsivo)
  {} % Nessuna indentazione
  {\bfseries} % Titolo in grassetto
  {.} % Punto dopo il titolo
  { } % Spazio dopo il titolo
  {} % Formato del titolo
% Per fare le matrici utilizzare \begin{pmatrix} a & b \\
%\usepackage{colortbl} % Per colorare celle e righe
\theoremstyle{nostyle}
\newtheorem{lemma}{Lemma}[section]
\newtheorem{theorem}{Teorema}[section]
\newtheorem{corollario}{Corollario}
\newtheorem{definition}{Definizione}[section]
\usepackage{enumerate}
\usepackage{tabularx}
\usepackage{nicematrix}
\usepackage{booktabs}
\usepackage{enumitem}
\usepackage{float}
\usepackage{xcolor}
\definecolor{verdeScuro}{RGB}{0, 100, 80}

\title{Basi di dati \\ \vspace{0.3cm}% chktex 8 
\large Maurizio Lenzerini}
\author{Nome autore} 
\affil{Università degli studi La Sapienza}
\affil{Ingegneria Informatica e Automatica - 3 anno} % chktex 8
\date{Gennaio - Dicembre 2026-27} % chktex 8



\begin{document}
\maketitle
\tableofcontents

\chapter{Modello Relazionale}
\section{Relazione matematica}
Siano $D_1,D_2,...,D_n$ n insiemi:\\
Il \textbf{prodotto cartesiano} $D_1 \times D_2 \times ... \times D_n$ è l'insieme di tutte le n-uple ordinate $(d_1,d_2,..,d_n)$ tali che $d_1 \in D_1, d_2 \in D_2 ,..., d_n \in D_n$. 
Una relazione matematica sugli insiemi $D_1,D_2,...,D_n$ è un \textbf{sottoinsieme del prodotto cartesiano} $D_1 \times D_2 \times ... \times D_n$.\\
Se R è una relaziome matematica sugli insiemi $D_1,D_2,...,D_n$ tali insiemi sono detti domini, tale relazione su $n$ domini dice che ha grado $n$. il numero di n-uple (righe) di una relazione è la cardinalità della relazione stessa.\\
Esempio:
\begin{itemize}

    \item $D_1 = \{\text{Conti}, \text{Totti}, \text{Maradona}\}$

    \item $D_2 = \{\text{Genoa}, \text{Roma}\}$

    \item prodotto cartesiano $D_1 \times D_2$
    
    \[
    \begin{array}{|l|l|}
    \hline
    \text{Conti}    & \text{Genoa} \\
    \text{Conti}    & \text{Roma}  \\
    \text{Totti}    & \text{Genoa} \\
    \text{Totti}    & \text{Roma}  \\
    \text{Maradona} & \text{Genoa} \\
    \text{Maradona} & \text{Roma}  \\
    \hline
    \end{array}
    \]

    \item una relazione:
    $\text{HaGiocato} \subseteq D_1 \times D_2$

    \[
    \begin{array}{|l|l|}
    \hline
    \text{Conti} & \text{Roma}  \\
    \text{Conti} & \text{Genoa} \\
    \text{Totti} & \text{Roma}  \\
    \hline
    \end{array}
    \]

\end{itemize}

\noindent La relazione matematica è quindi un insieme di ennuple o tuple ordinate su $D_1,D_2,...,D_n$, dove ogni tupla è una sequenza di valori della forma $(d_1,...,d_n)$.
Ciascuna tupla è ordinata nel senso che il valore del primo elemento della tupla proviene dal primo domonio, ma non è definito un ordinamento tra le tuple e sono tutte distinte.


\end{document}